\scanpage{267}
\mspage{781}

\unsure{mant|me} by the commutativity of the diagrams
\[
(107)\quad
\diagram{\begin{tikzcd}[ampersand replacement=\&, column sep=large,row sep=small]
H(A) \arrow[dd,"(100)"',"\sim"] \arrow[rd,"\substack{i_A\\\sim}"] \& \\
\& H(\unsure{\mathbb{S}}_x^{\sim},\unsure{\mathbb{D}}_x^{*\sim}) \\
H(A') \arrow[ru,"\sim","i_{A'}"'] \&
\end{tikzcd}}
\;.
\]

To verify the fact that the functor $\underline{\mathrm{Gr}}(H)$ of (102)
is \unsure{essentially} \unsure{constant},
I consider the two functors (independent of the choice of
$\unsure{\mathbb{S}}_x^{\sim},\unsure{\mathbb{D}}_x^{*\sim}$)
\[
(108)\quad
\left\{
\diagram{\begin{tikzcd}[ampersand replacement=\&, column sep=small,row sep=large]
A \arrow[r,mapsto] \& \unsure{\mathbb{S}}_{xA}^{\sim},
\& A \arrow[r,mapsto] \& \unsure{\mathbb{D}}_{x,A}^{*\sim} \\
\underline{\mathcal{A}}_x \arrow[r] \& \underline{\unsure{\mathrm{Revuni}}}(\unsure{\mathbb{S}}_x)
\& \underline{\mathcal{A}}_x \arrow[r] \& \underline{\unsure{\mathrm{Revuni}}}(\unsure{\mathbb{D}}_x^*)
\end{tikzcd}}
\right.
\;,
\]
which factor into \uline{equivalences}
\[
(109)\quad
\diagram{\begin{tikzcd}[ampersand replacement=\&, column sep=small]
\underline{\mathrm{Gr}}(\underline{\mathcal{A}}_x)
\arrow[r,"\substack{F\\\approx}"] \&
\underline{\unsure{\mathrm{Revuni}}}(\unsure{\mathbb{S}}_x),
\& \underline{\mathrm{Gr}}(\underline{\mathcal{A}}_x)
\arrow[r,"\substack{G\\\approx}"] \&
\underline{\unsure{\mathrm{Revuni}}}(\unsure{\mathbb{D}}_x^*)
\end{tikzcd}}
\]
This said, \inserted{using
$\unsure{\mathbb{S}}_x^{\sim},\unsure{\mathbb{D}}_x^{*\sim}$,}
we have for $A\in\mathrm{Ob}\,\underline{\mathcal{A}}_x=\mathcal{A}_x$ also,
isomorphisms functorial in $A$
\[
(110)\quad\left\{
\begin{array}{l}
\unsure{\mathbb{S}^{\sim}|H}(A)
\simeq\mathrm{Isom}_{\underline{\unsure{\mathrm{Revuni}}}(\unsure{\mathbb{S}}_x)}
\deleted{\unsure{\mathbb{S}}}
\left(\underbrace{F(A)}_{\substack{\Vert\\\unsure{\mathbb{S}}_{x,A}^{\sim}}},\unsure{\mathbb{S}}^{\sim}\right)
\\[1.8em]
\unsure{\mathbb{D}}_x^{*\sim}(\unsure{\mathrm{Ep}|E_p}A)
\simeq\mathrm{Isom}_{\underline{\unsure{\mathrm{Revuni}}}(\unsure{\mathbb{D}}_x^*)}
\left(\underbrace{G(A)}_{\substack{\Vert\\\unsure{\mathbb{D}}_{x,A}^{*\sim}}},\unsure{\mathbb{D}}_x^{*\sim}\right),
\end{array}\right.
\]
and by the universal property of $\underline{\mathrm{Gr}}(\underline{\mathcal{A}}_x)$,
these functorial isomorphisms extend over the category
$\underline{\mathrm{Gr}}(\underline{\mathcal{A}}_x)$.
Thus the situation is one where we have a commutative group $\unsure{T|\overline T}$ $(=T_x)$,
three connected $T$-groupoids $\mathcal{G},\mathcal{G}',\mathcal{G}''$

\margin{266}

\margin{\begin{tabular}{c}
266\\
\texttt{https://grothendieck.umontpellier.fr}\\
Université de Montpellier, UMR CNRS 5149, Case courrier 051, Place Eugène Bataillon\\
34095 Montpellier cedex 5 - France
\end{tabular}}

\scanpage{268}

(namely $\underline{\mathrm{Gr}}(\underline{\mathcal{A}}_x)$,
$\underline{\unsure{\mathrm{Revuni}}}(\unsure{\mathbb{S}}_x)$,
$\underline{\unsure{\mathrm{Revuni}}}(\unsure{\mathbb{D}}_x^*)$),
equivalences \deleted{\unsure{of}} of $T$-groupoids
\[
(111)\quad
\diagram{\begin{tikzcd}[ampersand replacement=\&, column sep=large]
\mathcal{G} \arrow[r,"\substack{F\\\approx}"] \& \mathcal{G}',
\& \mathcal{G} \arrow[r,"\substack{G\\\approx}"] \& \mathcal{G}''
\end{tikzcd}}
\]
and objects $x'\in\mathrm{Ob}\,\mathcal{G}'$,
$x''\in\mathrm{Ob}\,\mathcal{G}''$ (namely
$\unsure{\mathbb{S}}^{\sim}$ and $\unsure{\mathbb{D}}_x^{*\sim}$),
and if we consider the functors
\[
(112)\quad\left\{
\diagram{\begin{tikzcd}[ampersand replacement=\&, column sep=small,row sep=large]
\mathcal{G} \arrow[r,"f=f_{x'}"] \& \underline{\mathrm{Tors}}(T),
\& \mathcal{G} \arrow[r,"g=g_{x''}"] \& \underline{\mathrm{Tors}}(T) \\
x \arrow[r,mapsto] \& \mathrm{Hom}_{\mathcal{G}'}(F(x),x'),
\& x \arrow[r,mapsto] \& \mathrm{Hom}_{\mathcal{G}''}(G(x),x'')
\end{tikzcd}}
\right.
\]
and the corresponding functor
\[
(113)\quad\left\{
\diagram{\begin{tikzcd}[ampersand replacement=\&, column sep=large,row sep=large]
\mathcal{G} \arrow[r] \& \underline{\mathrm{Tors}}(T) \\
x \arrow[r,mapsto] \&
\mathrm{Hom}_{\underline{\mathrm{Tors}}(\unsure{T|G})}(f(x),g(x))
\end{tikzcd}}
\right.
\;,
\]
to prove that this last functor is essentially constant.
We have then, if \deleted{$F',G'$ denote}
\[
(114)\quad
\diagram{\begin{tikzcd}[ampersand replacement=\&, column sep=large]
F_1\deleted{\unsure{f}}:\mathcal{G}' \arrow[r,"\approx"] \& \mathcal{G},
\& \unsure{{}'G_1|G_1}:\mathcal{G}'' \arrow[r,"\approx"] \& \mathcal{G}
\end{tikzcd}}
\]
denote quasi-inverses of $F,G$, isomorphisms functorial in $x$
\[
(115)\quad\left\{
\begin{array}{l}
f(x)\overset{\mathrm{def}}{=}\mathrm{Hom}_{\mathcal{G}'}(F(x),x')
\simeq\mathrm{Hom}_{\mathcal{G}}(x,F_1(x')) \\
g(x)\overset{\mathrm{def}}{=}\mathrm{Hom}_{\mathcal{G}''}(G(x),x'')
\simeq\mathrm{Hom}_{\mathcal{G}}(x,F_2(x''))
\end{array}\right.
\]
whence \unsure{follows|follow} isomorphisms of $T$-torsors
\[
(116)\quad
\unsure{\mathrm{Isom}}(f(x),g(x))
\simeq\unsure{\mathrm{Isom}}(F_1(x'),F_2(x''))
\]
whence the assertion.

\margin{267}

\margin{\begin{tabular}{c}
267\\
\texttt{https://grothendieck.umontpellier.fr}\\
Université de Montpellier, UMR CNRS 5149, Case courrier 051, Place Eugène Bataillon\\
34095 Montpellier cedex 5 - France
\end{tabular}}

\scanpage{269}
\mspage{782}

\uline{Definition} We say that we have a \uline{correlation}
between a universal covering $\unsure{\mathbb{S}}_x^{\sim}$ of
$\unsure{\mathbb{S}}_x$, and another $\unsure{\mathbb{D}}_x^{*\sim}$
of $\unsure{\mathbb{D}}_x^*$, when we have given ourselves an element of
$H(\unsure{\mathbb{S}}_x^{\sim},\unsure{\mathbb{D}}_x^{*\sim})$.

Thus, if $A\in\mathcal{A}_x$, the datum of a correlation amounts to the
datum of an element of $H(A)$, i.e. of an isomorphism
\deleted{\unsure{of}} between the $T_x$-torsors
$\unsure{\mathbb{S}}_x^{\sim}(A)$ and
$\unsure{\mathbb{D}}_x^{*\sim}(\unsure{\mathrm{Ep}|E_p}(A))$.
This identification is compatible with the inclusion relations
between arcs $A\subset A'$.
\deleted{\unsure{D}} The correlations between
$\unsure{\mathbb{S}}_x^{\unsure{*|\sim}}$ and
$\unsure{\mathbb{D}}_x^{*\sim}$ form a torsor \unsure{under|over} $T_x$.
\deleted{For $\unsure{\mathbb{S}}_x^{\sim}$ fixed}

Let $\underline{\mathrm{Corr}}_x(\unsure{\mathfrak{X}|X})
=\underline{\mathrm{Corr}}_x(\deleted{\unsure{\mathbb{S}}}\unsure{X})$
be the category of \deleted{\unsure{pairs}}\unsure{triples}
$(\widetilde{\unsure{\mathbb{S}}}_x,
\widetilde{\unsure{\mathbb{D}}}_x^*,h)$ of correlated universal
coverings of $\unsure{\mathbb{S}}_x$ and of $\unsure{\mathbb{D}}_x^*$,
we have \deleted{\unsure{equi|equivalences}} functors
\[
\diagram{\begin{tikzcd}[ampersand replacement=\&, column sep=large,row sep=0.2ex]
\& \unsure{\mathbb{S}}_x^{\unsure{*\sim|\sim}} \\
(\widetilde{\unsure{\mathbb{S}}}_x,\unsure{\mathbb{D}}_x^{*\sim},h)
\arrow[ru] \arrow[rd] \& \\
\& \unsure{\mathbb{D}}_x^{*\sim}
\end{tikzcd}}
\]
\[
(117)\quad
\diagram{\begin{tikzcd}[ampersand replacement=\&, column sep=large,row sep=large]
\& \underline{\mathrm{Corr}}_x(X) \arrow[ld,"\approx"'] \arrow[rd,"\approx"] \& \\
\underline{\unsure{\mathrm{Revuni}}}(\unsure{\mathbb{S}}_x)
\& \& \underline{\unsure{\mathrm{Revuni}}}
(\deleted{\unsure{\mathbb{S}}}\unsure{\mathbb{D}}_x^*)
\end{tikzcd}}
\;,
\]
\unsure{which} are equivalences of categories.

\margin{268}

\margin{\begin{tabular}{c}
268\\
\texttt{https://grothendieck.umontpellier.fr}\\
Université de Montpellier, UMR CNRS 5149, Case courrier 051, Place Eugène Bataillon\\
34095 Montpellier cedex 5 - France
\end{tabular}}

\scanpage{270}

All these constructions have been carried out intrinsically carefully
enough to \unsure{pass} to isomorphisms of germs of differentiable
surfaces
$\diagram{\begin{tikzcd}[ampersand replacement=\&, column sep=large]
(X,x) \arrow[r,"\substack{u\\\sim}"] \& (X',x')
\end{tikzcd}}$.
Suppose chosen, in addition to
$(\unsure{\mathbb{S}}_x^{\sim},\unsure{\mathbb{D}}_x^{*\sim},h_x)$
\unsure{relative|relative} to $(X,x)$, also
$(\unsure{\mathbb{S}}_{x'}^{\sim},\unsure{a}\,
\unsure{\mathbb{D}}_{x'}^{*\sim},h'_{x'})$ \unsure{relative|relative} to $(X',x')$.
\unsure{Now} the isomorphism of the germs being given, it therefore induces
\[
(118)\quad\left\{
\diagram{\begin{array}{lll}
u_{\unsure{\mathbb{S}}}:&
\unsure{\mathbb{S}}_x\simeq\unsure{\mathbb{S}}_{x'}
&\text{isomorphism of topological spaces}\\[0.7em]
u_{\unsure{\mathbb{D}}^*}:&
\unsure{\mathbb{D}}_x^*\xrightarrow{\sim}\unsure{\mathbb{D}}_{x'}^*
&\text{isomorphism of germs of spaces}
\end{array}}
\right.
\]
\[
\diagram{\begin{array}{lll}
u_T:&T_x\simeq T_{x'}
&\begin{array}{l}
\text{isomorphism of the }\mathbb{Z}\text{-modules }\inserted{\text{topological}}\\
\unsure{\text{free}}\text{ of rank }1
\end{array}
\end{array}}
\]
Let us then consider the two sets
\[
(119)\quad\left\{
\begin{array}{l}
F(u)=\text{set of }\unsure{\text{isomorphisms}}\quad
\unsure{\mathbb{S}}_x^{\sim}\simeq\unsure{\mathbb{S}}_{x'}^{\sim}
\text{ over }\inserted{u_{\unsure{\mathbb{S}}}}
\\[0.6em]
\hfill\deleted{\unsure{u_{\mathbb{S}}}}
\\[0.6em]
G(u)=\text{set of isomorphisms}\quad
\unsure{\mathbb{D}}_x^{*\sim}\simeq\unsure{\mathbb{D}}_{x'}^{*\sim}
\text{ over }\inserted{u_{\unsure{\mathbb{D}}^*}},
\end{array}\right.
\]
they are both $T_x$-torsors.
This being so, I say that we have a canonical isomorphism of $T_x$-torsors
\[
(120)\quad
\diagram{\begin{tikzcd}[ampersand replacement=\&, column sep=large]
F(u) \arrow[r,"\sim"] \& G(u)
\end{tikzcd}}
\]
--- i.e. ``it comes to the $=$ to lift $u_{\unsure{\mathbb{S}}}$ to
$u_{\widetilde{\unsure{\mathbb{S}}}}:
\unsure{\mathbb{S}}_x^{\sim}\simeq S_{x'}^{\sim}$,
or to lift $u_{\unsure{\mathbb{D}}^*}$ to
\deleted{$\unsure{u_{\widetilde{\mathbb{S}}}}$}
$u_{\unsure{\mathbb{D}}^{*\sim}}:
\unsure{\mathbb{D}}_x^{*\sim}\unsure{\longrightarrow|-}
\unsure{\mathbb{D}}_{x'}^{*\sim}$''.
To see this, we note that by
\deleted{functoriality of the \unsure{constructions}}
\inserted{transport of structure} the

\margin{269}

\margin{\begin{tabular}{c}
269\\
\texttt{https://grothendieck.umontpellier.fr}\\
Université de Montpellier, UMR CNRS 5149, Case courrier 051, Place Eugène Bataillon\\
34095 Montpellier cedex 5 - France
\end{tabular}}

\scanpage{271}
\mspage{783}

inverse images under $u_{\unsure{\mathbb{S}|\mathcal{S}}}$ and under $u_{\unsure{\mathbb{D}|D}^{*}}$ of
$\unsure{\mathbb{S}|\mathcal{S}}_{x'}^{\sim}$ and of
$\unsure{\mathbb{D}|D}_{x'}^{*\sim}$,
\unsure{that is|that is} $\unsure{\mathbb{S}|\mathcal{S}}_x^{\prime\sim}$, $\unsure{\mathbb{D}|D}_x^{*\prime\sim}$,
are equipped with a correlation, which I denote by $h'_x$.
This posited, \deleted{\unsure{the existence} of an isomorphism
$F(u)\simeq G(u)$} canonical isomorphisms of $T_x$-torsors
\[
\begin{array}{rl}
(121)&F(u)\simeq\operatorname{Isom}_{\unsure{\mathbb{S}|\mathcal{S}}_x}
  (\unsure{\mathbb{S}|\mathcal{S}}_x^{\sim},\unsure{\mathbb{S}|\mathcal{S}}_x^{\prime\sim})\\[3pt]
&G(u)\simeq\operatorname{Isom}_{\unsure{\mathbb{D}|D}_x^{*}}
  (\unsure{\mathbb{D}|D}_x^{*\sim},\unsure{\mathbb{D}|D}_x^{*\prime\sim})
\end{array}
\]
Now by the diagram of equivalences (117) we do have isomorphisms of
$T_x$-torsors
\[
\begin{gathered}(122)\\[3pt]
\diagram{\begin{tikzcd}[ampersand replacement=\&, column sep=small,row sep=large,column sep=9pt,row sep=14pt,execute at begin picture={\fontsize{10}{12}\selectfont}]
\& {\begin{gathered}\operatorname{Isom}((\unsure{\mathbb{S}|\mathcal{S}}_x^{\sim},\unsure{\mathbb{D}|D}_x^{*\sim},h_x),\\[-1pt]
(\unsure{\mathbb{S}|\mathcal{S}}_x^{\prime\sim},\unsure{\mathbb{D}|D}_x^{*\prime\sim},h'_x))\end{gathered}}
\arrow[dl,"\sim"']\arrow[dr,"\sim"] \& \\
{\begin{gathered}F(u)=\\[-1pt]\operatorname{Isom}(\unsure{\mathbb{S}|\mathcal{S}}_x^{\sim},\unsure{\mathbb{S}|\mathcal{S}}_x^{\prime\sim}\unsure{]|)}\end{gathered}}
\&\&{\begin{gathered}G(u)=\\[-1pt]\unsure{\operatorname{Isom}}(\unsure{\mathbb{D}|D}_x^{*\sim},\unsure{\mathbb{D}|D}_x^{*\prime\sim})\end{gathered}}
\end{tikzcd}}
\quad ,
\end{gathered}\]
whence the canonical bijection \unsure{sought} (120) ---
\deleted{\unsure{of the}} which gives in particular (96).

\deleted{\unsure{The}} The canonical bijection (120) has the following
$\pm$ obvious property, which characterizes it:
\[
(123)\quad\left\{\vcenter{\hbox{\begin{minipage}{0.82\linewidth}
\deleted{\unsure{We}} Let $u_{\unsure{\mathbb{S}|\mathcal{S}}^{\sim}}$ and
$u_{\unsure{\mathbb{D}|D}^{*\sim}}$ be
\deleted{\unsure{extensions}}\inserted{liftings} of $u$
\deleted{\unsure{$u_{\mathbb{S}}$}}\deleted{\unsure{associated}}
$\unsure{\mathbb{S}|\mathcal{S}}_x^{\sim}\simeq\unsure{\mathbb{S}|\mathcal{S}}_{x'}^{\sim}$,
$\unsure{\mathbb{D}|D}_x^{*\sim}\longrightarrow\unsure{\mathbb{D}|D}_{x'}^{*\sim}$,
associated by (120), let $A\in\mathcal{A}_x$,
$A'\in\mathcal{A}'_x$ its image under
$u_{\unsure{\mathbb{S}|\mathcal{S}}}$, let $\alpha$ be a section of $\unsure{\mathbb{S}|\mathcal{S}}_x^{\sim}$ over $A$,
$\beta$ the section of $\unsure{\mathbb{D}|D}_x^{*\sim}$ over
$\unsure{\mathrm{Ep}|E_p}(A)$ \unsure{associated|associated} by the correlation
$h_x$, let $\alpha'$ be the section of $\unsure{\mathbb{S}|\mathcal{S}}_{x'}^{\sim}$ over $A'$
deduced from $\alpha$ by applying $u_{\unsure{\mathbb{S}|\mathcal{S}}^{\sim}}$,
$\beta'$ the section of $\unsure{\mathbb{D}|D}_{x'}^{*\sim}$ over
$\mathrm{Ep}(A')$ deduced from $\beta$ by applying
$u_{\unsure{\mathbb{D}|D}^{*\sim}}$, then $\alpha'$ and $\beta'$ are
\unsure{associated} by the correlation $h'$.
\end{minipage}}}\right.
\]
\margin{270}

\margin{\begin{tabular}{c}
270\\
\texttt{https://grothendieck.umontpellier.fr}\\
Université de Montpellier, UMR CNRS 5149, Case courrier 051, Place Eugène Bataillon\\
34095 Montpellier cedex 5 - France
\end{tabular}}

\scanpage{272}

This \deleted{\unsure{characterizes}} \unsure{(}for $A,\alpha$ fixed)
characterizes the map \deleted{bijection} (120), because for
$u_{\unsure{\mathbb{S}|\mathcal{S}}^{\sim}}$ fixed, it determines $u_{\unsure{\mathbb{D}|D}^{*\sim}}$ as follows:
we have $\beta$ by $h_x$, $\alpha'$ \unsure{deduced|deduced} from $\alpha$ by
$u_{\unsure{\mathbb{S}|\mathcal{S}}}$, $\beta'$ \unsure{deduced|deduced} from $\alpha'$ by $h_{x'}$,
now $\beta'$ must be equal to the section deduced from $\beta$ by applying
$u_{\unsure{\mathbb{D}|D}^{*\sim}}:\unsure{\mathbb{D}|D}_x^{*\sim}
\unsure{\longrightarrow|-}\unsure{\mathbb{D}|D}_{x'}^{*\sim}$,
and this indeed determines $u_{\unsure{\mathbb{D}|D}^{*\sim}}$ uniquely ...

\uline{N.B.} We can use this construction, for $A,\alpha$ fixed,
to define the map (120), we find that it is a
$T_x$-morphism (hence an isomorphism) and,
\deleted{\unsure{that it|since it}}\inserted{at the \unsure{same|=}}
time, that it does not depend on the choice of $\alpha$.
But a priori it depends \unsure{perhaps} on the choice of $A$, and the point which
requires verification is that this bijection in fact does not depend on the choice
of $A$. We know indeed that for $A\subset A'$, it is the
\unsure{same|=} for $A$ and $A'\unsure{;|,}$
whence at once the result, for the equivalence relation generated in
$\mathcal{A}_x$ by the inclusion relation is the indiscrete relation .... .

\margin{271}

\margin{\begin{tabular}{c}
271\\
\texttt{https://grothendieck.umontpellier.fr}\\
Université de Montpellier, UMR CNRS 5149, Case courrier 051, Place Eugène Bataillon\\
34095 Montpellier cedex 5 - France
\end{tabular}}

\scanpage{273}
\mspage{784}

F) \uline{Application} : \uline{differential interpretation of the special
Teichmüller groups}.

Let $X$ be a topological space, equipped with a subset $I$ that is discrete
\inserted{finite}, such that $X$ is a surface in a neighborhood of $I$ ---
moreover we give ourselves a differentiable structure \inserted{$(C_1)$}
in a neighborhood \unsure{of the|of} $I$. Let $G$ \deleted{\illegible{1}} be a
topological group, operating continuously on $(X,I)$, and respecting the
differentiable structure of $X$ in a neighborhood of $S$. \unsure{Let|Let}
\[
(124)\quad\left\{\begin{array}{l}
\unsure{\mathbb{S}|\mathcal{S}}_{X,I}\ \unsure{\text{or}|\sim}\ \unsure{\mathbb{S}|\mathcal{S}}_I
=\displaystyle\coprod_{x\in I}
\underbrace{\unsure{\mathbb{S}|\mathcal{S}}_{X,x}}_{\text{or }\unsure{\mathbb{S}|\mathcal{S}}_x}\ ,\\[8pt]
\unsure{\mathbb{D}|D}_{X,I}^{*}\ \unsure{\text{or}|\sim}\ \unsure{\mathbb{D}|D}_{I}^{*}
=\displaystyle\coprod_{x\in I}
\underbrace{\unsure{\mathbb{D}|D}_{X,x}^{*}}_{\text{or }\unsure{\mathbb{D}|D}_x^{*}}
\end{array}\right.
\]
we assume that the operation of $G$ on $\unsure{\mathbb{S}|\mathcal{S}}_S$ induced by
the action of $G$ is also continuous.

Let, for $\forall\ x\in I$, $\unsure{\mathbb{D}|D}_x^{*\sim}$ be a universal
covering of $\unsure{\mathbb{D}|D}_x^{*}$, and put
\[
(125)\quad
\unsure{\mathbb{D}|D}_{X,I}^{*\sim}\ \unsure{\text{or}|-}\ \unsure{\mathbb{D}|D}_I^{*\sim}
=\coprod_{x\in S}\unsure{\mathbb{D}|D}_x^{*\sim}\ .
\]
Let us for the moment forget the differentiable \unsure{data},
and consider the group
\[
(126)\quad
\widetilde{G}=\left\{u,\widetilde{u}\ \middle|\
\begin{array}{l}
u\in G,\ \widetilde{u}\text{ a lifting of}\\
u_{\unsure{\mathbb{D}|D}_I^{*}}\text{ to an \unsure{automorphism} }
 u_{\unsure{\mathbb{D}|D}_I^{*\sim}}\\
\text{of }\unsure{\mathbb{D}|D}_I^{*\sim}
\end{array}\right\}
\]
We have a homomorphism $(u,\widetilde{u})\longmapsto\widetilde{u}$,
$\widetilde{G}\longrightarrow G$, which is

\margin{272}

\margin{\begin{tabular}{c}
272\\
\texttt{https://grothendieck.umontpellier.fr}\\
Université de Montpellier, UMR CNRS 5149, Case courrier 051, Place Eugène Bataillon\\
34095 Montpellier cedex 5 - France
\end{tabular}}

\scanpage{274}

an epimorphism, and, putting
\[
(127)\quad T_I=\prod_{x\in I}T_x\ ,
\]
we have an exact sequence
\[
(128)\quad
\diagram{\begin{tikzcd}[ampersand replacement=\&, column sep=small]
1\arrow[r] \& T_I\arrow[r,thick] \& \widetilde{G}\arrow[r] \& G\arrow[r] \& 1
\end{tikzcd}}\ .
\]
I would like to introduce on $\widetilde{G}$ a topology such that (128)
becomes an exact sequence of topological groups.
Replacing $G$ if necessary by the open subgroup $G'$ of finite index
kernel of the canonical homomorphism
\[
(129)\quad
\diagram{\begin{tikzcd}[ampersand replacement=\&]
G\arrow[r] \& \unsure{\mathfrak{S}|\mathcal{S}}_I
\end{tikzcd}}\ ,
\]
hence $\widetilde{G}$ by the inverse image $\widetilde{G}'$ of $G'$,
we are reduced to the case where $G$ \uline{operates trivially on} $I$,
\deleted{and If need be \illegible{1}}
\inserted{Then $\widetilde{G}$ is expressed}
\deleted{further to a smaller open \unsure{subgroup}
(of index dividing $2^{\operatorname{card}I}$)
\illegible{1}}
\deleted{\unsure{= that}}
as the fiber product over $G$ of the family, indexed by $I$,
of the extensions
\[
(130)\quad
\diagram{\begin{tikzcd}[ampersand replacement=\&, column sep=small]
1\arrow[r] \& T_x\arrow[r] \& \widetilde{G}_x\arrow[r] \& G\arrow[r] \& 1
\end{tikzcd}}
\]
corresponding to the cases $I_x\overset{\mathrm{def}}{=}\{x\}$.
Thus we are reduced to the case where $\operatorname{card}\unsure{I|S}=1$.

\margin{273}

\margin{\begin{tabular}{c}
273\\
\texttt{https://grothendieck.umontpellier.fr}\\
Université de Montpellier, UMR CNRS 5149, Case courrier 051, Place Eugène Bataillon\\
34095 Montpellier cedex 5 - France
\end{tabular}}

\scanpage{275}
\mspage{785}

Let \deleted{$\unsure{\Delta_x|\mathbb D_x}$}
\[
(131)\qquad
\unsure{\{}\Delta_{x,0}\deleted{\unsure{\}_{\alpha\in\mathbf N}}}
\]
be \unsure{a|a} \deleted{\unsure{decreasing sequence}} of \unsure{neighborhoods|neighborhood} of $x$ in $X$ which \unsure{is|are} \deleted{\unsure{some}} \inserted{an} \unsure{open disks|open disk}, \deleted{and that \unsure{form} a \unsure{basis} of the filter of neighborhoods of $x$ in $X$, and let}
We can \unsure{then} find a universal covering $\Delta_{x,0}^{*\sim}$ of $\Delta_{x,0}^{*}=\Delta_{x,0}\setminus\{x\}$ which \unsure{induces|induces} $\unsure{\mathbb D|D}_x^{*\sim}$ on $\unsure{\mathbb D|D}_x^{*}$. By the continuity hypothesis, there exists a neighborhood $U$ of $1$ in $G$, and a neighborhood $\Delta_{x,1}$ of $x$, such that we have
\[
(132)\qquad
\diagram{\begin{tikzcd}[ampersand replacement=\&]
U\times\Delta_{x,1} \arrow[r] \arrow[d, phantom, "\cap" description] \& \Delta_{x,0} \arrow[d, phantom, "\cap" description] \\
G\times X \arrow[r] \& X
\end{tikzcd}},
\]
and we \unsure{can} suppose that $\Delta_{x,1}$ is an \unsure{open disk|open disks}. This taken; for fixed $u\in\unsure{U}$, giving the lifting $\widetilde u$ of $u_{\unsure{\mathbb D|D}_x}$ to $\unsure{\mathbb D|D}_x^{\sim}$ is equivalent to giving a lifting to \deleted{$\unsure{\Delta|\mathbb D}$}
\[
(133)\qquad
\diagram{\begin{tikzcd}[ampersand replacement=\&]
\widetilde\Delta_{x,1} \arrow[r, "\widetilde u_{0,1}"] \& \Delta_{x,0}^{\sim}
\end{tikzcd}}
\]
compatible with
\[
u\mid\Delta_{x,1}^{*}:\quad
\diagram{\begin{tikzcd}[ampersand replacement=\&]
\Delta_{x,\unsure{1|0}}^{*}\arrow[r] \& \Delta_{x,0}^{*}
\end{tikzcd}}.
\]
That said, in the inverse image $\widetilde U$ in $\widetilde G$ of the neighborhood $U$ of $1$ in $G$, we are going to

\margin{274}

\margin{\begin{tabular}{c}
274\\
\texttt{https://grothendieck.umontpellier.fr}\\
Université de Montpellier, UMR CNRS 5149, Case courrier 051, Place Eugène Bataillon\\
34095 Montpellier cedex 5 - France
\end{tabular}}

\scanpage{276}

introduce the topology \inserted{the coarsest making continuous} \deleted{of \unsure{compact convergence} \unsure{of the|of} $\widetilde u$ and of $u$} the maps
\[
(134)\quad
\left\{\begin{array}{ll}
\diagram{\begin{tikzcd}[ampersand replacement=\&]
\unsure{p_1|p_i}:\widetilde U \arrow[r] \& U
\end{tikzcd}}
& \text{\unsure{induced|induced} by }\diagram{\begin{tikzcd}[ampersand replacement=\&]\widetilde G\arrow[r]\&G\end{tikzcd}}
\\[1.2ex]
\diagram{\begin{tikzcd}[ampersand replacement=\&]
q:\widetilde U \arrow[r, phantom, "\unsure{\Longrightarrow|\longrightarrow}" description] \arrow[r, phantom, "\deleted{\unsure{\longrightarrow}}"', shift right=2ex] \&[35pt] \operatorname{Cont}(\Delta_{x,1}^{*\sim},\Delta_{x,0}^{*\sim})
\end{tikzcd}}
& \diagram{\begin{tikzcd}[ampersand replacement=\&]
(u,\widetilde u) \arrow[r, mapsto] \& \deleted{\unsure{(\widetilde u)_{0,1}}}\inserted{\widetilde u_{0,1}}
\end{tikzcd}}
\\[-.2ex]
& \diagram{\begin{tikzcd}[ampersand replacement=\&,nodes in empty cells]
\deleted{\unsure{(u,\widetilde u)}} \arrow[r, phantom, "\deleted{\mapsto}" description] \& \deleted{(\widetilde u^{-1})_{0,1}}
\end{tikzcd}}
\\[1.2ex]
q'=q\circ\sigma,\quad\text{where }\sigma(u,\widetilde u)=(u,\widetilde u)^{-1}: & \diagram{\begin{tikzcd}[ampersand replacement=\&]
\widetilde U\arrow[r]\&\operatorname{Cont}(\Delta_{x,1}^{*\sim},\Delta_{x,0}^{*\sim})
\end{tikzcd}}
\end{array}\right.
\]

\unsure{We shall have|I shall have} taken care to take $\unsure{U|\widetilde U}$ stable under $u\mapsto u^{-1}$, which implies that the \unsure{same|=} holds for $\widetilde U$.

We thus find, with the help of this topology on $\widetilde U$, \deleted{\unsure{a} \illegible{1} \unsure{a filter}} the filter of neighborhoods of $1_{\widetilde G}$ in $\widetilde U$, hence also its \unsure{image} filter in $\widetilde G$. We must check that the latter is indeed the filter of neighborhoods of $1$ for a topological group topology, and that for this topology, $\widetilde G\to G$ is a covering of $G$. \deleted{\unsure{I suppose}}\inserted{I admit} that everything works.

When $X$ is \inserted{a surface} compact and orientable, \inserted{connected}, and \deleted{that} if $\unsure{T|\overline T}$ denotes its module of

\margin{275}

\margin{\begin{tabular}{c}
275\\
\texttt{https://grothendieck.umontpellier.fr}\\
Université de Montpellier, UMR CNRS 5149, Case courrier 051, Place Eugène Bataillon\\
34095 Montpellier cedex 5 - France
\end{tabular}}

\scanpage{277}
\mspage{786}

orientations, then all the $T_x$ are canonically isomorphic to $\unsure{T|\overline T}$
\[
(135)\qquad T_x\simeq\unsure{T|\overline T}\qquad\forall x\in I,
\]
and the exact sequence (128) takes the form
\[
(136)\qquad
\diagram{\begin{tikzcd}[ampersand replacement=\&]
1\arrow[r, dash]\&\unsure{T^I|\overline T^I}\arrow[r]\&\widetilde G\arrow[r]\&G\arrow[r]\&1
\end{tikzcd}}
\]
The ``universal'' case is the one where
\[
(137)\qquad G=\mathcal A=\operatorname{Aut\,top}(X,S)
\]
(with the topology of uniform convergence of the $u$ and of the $u^{-1}$), Then in the \uline{\unsure{anabelian}} cases (the type of $(X,S)$ being $g,\nu$, and
\[
(138)\qquad\begin{array}{l}
2g-2+\nu>0\quad\text{i.e.}\quad 2g+\nu\geq3\\[.5ex]
\text{i.e.}\quad(g,\nu)\ne(0,0),(0,1),(0,2),(1,0)
\end{array}
\]
\unsure{)|,} then we know that $\mathcal A^0$ is \unsure{1-connected} (\unsure{is|and} \unsure{even|=} contractible) hence \unsure{every|every} \deleted{\unsure{covering|covering of}}\inserted{extension} of $\mathcal A^0$ by a discrete group is canonically split, and the extension $\widetilde{\mathcal A}$ of $G$ then comes from an extension $\widetilde{\unsure{\mathcal T|\tau}}$
\[
(139)\qquad
\diagram{\begin{tikzcd}[ampersand replacement=\&]
1\arrow[r]\&\unsure{\mathbb Z^I|T^I}\arrow[r, dash]\&\widetilde{\unsure{\mathcal T|\tau}}\arrow[r]\&\unsure{\mathcal T|\tau}\arrow[r]\&1
\end{tikzcd}}
\]
of the Teichmüller group
\[
(140)\qquad
\unsure{\mathcal T|\tau}=\unsure{\mathcal T|\tau}(X,S)=\mathcal A/\mathcal A^0=\pi_0\bigl(\underbrace{\uline{\operatorname{Aut\,top}}(X,S)}_{\unsure{A|\mathcal A}}\bigr)
\]

\margin{276}

\margin{\begin{tabular}{c}
276\\
\texttt{https://grothendieck.umontpellier.fr}\\
Université de Montpellier, UMR CNRS 5149, Case courrier 051, Place Eugène Bataillon\\
34095 Montpellier cedex 5 - France
\end{tabular}}

\scanpage{278}

\unsure{which|that} thus takes the name of \uline{special
Teichmüller group} of $(X,S)$. For
arbitrary $G$ the action of $G$ on
$(X,S)$ \unsure{amounts} to giving a homomorphism
of topological groups
\diagram{\begin{tikzcd}[ampersand replacement=\&]
G \arrow[r] \& \unsure{\mathcal A|A}
\end{tikzcd}}
\unsure{fitting} into a diagram of
topological groups
\[
\tag{141}
\diagram{\begin{tikzcd}[ampersand replacement=\&]
G \arrow[r] \arrow[d] \& \unsure{\mathcal A|A} \arrow[d] \\
\pi_0(G) \arrow[r] \& \unsure{\mathcal T|\tau}=\pi_0(\unsure{\mathcal A|A})
\end{tikzcd}},
\]
and we then find, by functoriality,
that the $\widetilde G$ constructed previously
is the inverse image of \inserted{the extension}
$\widetilde{\unsure{\mathcal T|\tau}}=\operatorname{Sp}\unsure{\mathcal T|\tau}$
$\unsure{a|\alpha}$ of $\unsure{\mathcal T|\tau}$
by $\mathbb Z^I$, via $G\longrightarrow\unsure{\mathcal T|\tau}$.

\unsure{These} preliminaries being \unsure{taken},
let us return to the situation where $X$
admits a differentiable structure $C_1$ in a neighborhood
of $\unsure{I|S}$, and let for $\forall\,x\in\unsure{I|S}$,

\margin{277}

\margin{\begin{tabular}{c}
277\\
\texttt{https://grothendieck.umontpellier.fr}\\
Université de Montpellier, UMR CNRS 5149, Case courrier 051, Place Eugène Bataillon\\
34095 Montpellier cedex 5 - France
\end{tabular}}

\scanpage{279}
\mspage{787}

$\widetilde{\unsure{\mathbb S|\mathcal S}}_x$
be ``\unsure{the}'' universal covering of $\unsure{\mathbb S|\mathcal S}_x$ \unsure{correlated} with the
universal covering $\unsure{\mathbb D|D}_x^{*\sim}$ of $\unsure{\mathbb D|D}_x^*$. Put
\[
\tag{142}
\widetilde{\unsure{\mathbb S|\mathcal S}}_I
=\coprod_{x\in I}\widetilde{\unsure{\mathbb S|\mathcal S}}_x,
\]
and note that we have, for $\forall\,u\in G$,
a one-to-one correspondence between
the set of liftings $\widetilde u_{\unsure{\mathbb D|D}_I^*}$ of $u_{\unsure{\mathbb D|D}_I^*}$
to automorphisms of $\unsure{\mathbb D|D}_I^{*\sim}$, and the set of
liftings $\widetilde u_{\unsure{\mathbb S|\mathcal S}_I}$ of
$u_{\unsure{\mathbb S|\mathcal S}}$ to automorphisms of
$\widetilde{\unsure{\mathbb S|\mathcal S}}_I$,
which is an isomorphism of $T_I$-torsors.
This correspondence is compatible with
compositions, and we thus find
a canonical isomorphism
\[
\tag{143}
\widetilde G\simeq\text{group of the }
\left\{u,\widetilde u_{\unsure{\mathbb S|\mathcal S}}\;\middle|\;
\begin{array}{l}
u\in G,\ \widetilde u_{\unsure{\mathbb S|\mathcal S}_I}\ \text{lifting}\\
\text{of }u_{\unsure{\mathbb S|\mathcal S}_I}\ \text{to an}\\
\text{automorphism of }\unsure{\mathbb S|\mathcal S}_I^{\sim}
\end{array}\right\}
\]

We check that this is a \uline{topological}
\uline{isomorphism} by putting on the second
\deleted{\unsure{met}} the coarsest topology making
continuous the maps
\[
\tag{144}
\begin{array}{l}
\diagram{\begin{tikzcd}[ampersand replacement=\&, column sep=large]
\unsure{p_1^{a}|p_1}:\widetilde G \arrow[r] \& G
\end{tikzcd}}\\
\diagram{\begin{tikzcd}[ampersand replacement=\&, column sep=large]
q_1:\widetilde G \arrow[r] \&
\operatorname{Aut\,top}(\widetilde{\unsure{\mathbb S|\mathcal S}}_I)
\end{tikzcd}}
\qquad (u,\widetilde u_{\unsure{\mathbb S|\mathcal S}})
\longmapsto\widetilde u_{\unsure{\mathbb S|\mathcal S}}\\
\diagram{\begin{tikzcd}[ampersand replacement=\&, column sep=large]
q'_1=q_1\circ\sigma:\widetilde G \arrow[r] \&
\operatorname{Aut\,top}(\widetilde{\unsure{\mathbb S|\mathcal S}}_I)
\end{tikzcd}}
\end{array}
\]
\margin{\unsure{trans|tras}}

\margin{278}

\margin{\begin{tabular}{c}
278\\
\texttt{https://grothendieck.umontpellier.fr}\\
Université de Montpellier, UMR CNRS 5149, Case courrier 051, Place Eugène Bataillon\\
34095 Montpellier cedex 5 - France
\end{tabular}}

\scanpage{280}

where $\operatorname{Aut\,top}(\widetilde{\unsure{\mathbb S|\mathcal S}}_I)$ is equipped with
compact convergence.

In the case where $X$ is \inserted{a surface} compact and differentiable
and orientable, we thus find an interpretation
of the preceding groups,
\deleted{\unsure{constructed|explicit}, \unsure{by}}
\unsure{even|way} differential,
\inserted{in \unsure{so far} \unsure{as an extension|as an extension} of $\unsure{G|\widetilde G}$ by $\unsure{T_I|\overline T_I}$,}
as the inverse
image of a special Teichmüller
group $\unsure{S\mathcal T|\operatorname{Sp}\mathcal T}\longrightarrow\unsure{\mathcal T|\tau}$,
by the canonical homomorphism
$G\longrightarrow\unsure{\mathcal T|\tau}$ of (141).

\margin{279}

\margin{\begin{tabular}{c}
279\\
\texttt{https://grothendieck.umontpellier.fr}\\
Université de Montpellier, UMR CNRS 5149, Case courrier 051, Place Eugène Bataillon\\
34095 Montpellier cedex 5 - France
\end{tabular}}
